\begin{theorem}[The literal method is not graph-uniformly accelerated]
\label{thm:direct-three-vertex-lower-bound}
Fix $B=5/2$ and $\epsppr=10^{-2}$.  Let $P_3$ be the three-vertex path,
seeded at its degree-two center, and run the literal top-$1/32$ spreading
phase with optimal SOR.  For every $0<\alpha\leq10^{-4}$,
\begin{equation}
 \Work_{\rm spread}\geq\frac1{1408}\alpha^{-3/2}.
 \label{eq:direct-three-vertex-lower-bound}
\end{equation}
Here $\mathcal V_{\rm exp}=4$.  Therefore a graph-uniform bound
$\widetilde O(\mathcal V_{\rm exp}/\sqrt\alpha)$ for the literal two-rung
method is false, even on one fixed graph and before the exact rung begins.
\end{theorem}

\begin{proof}
Put
\begin{equation*}
 h:=\sqrt\alpha,
 \quad \lambda:=\frac{1-h}{1+h},
 \quad \eta:=\frac{\lambda}{1+\lambda^2},
 \quad e:=2\eta,
 \quad \zeta:=\frac\alpha2.
\end{equation*}
At a radial state let $X=y_{\rm root}/\zeta$ and let
$Y=y_{\rm leaf}/\zeta$ be the common normalized residual at either leaf.
The scaled spreading gate is
\begin{equation*}
 \beta:=\frac{B\alpha\epsppr}{\zeta}
 =2B\epsppr=\frac1{20}.
\end{equation*}
There are at most three live coordinates, so every literal batch has size
$\lceil|\mathcal F|/32\rceil=1$.  A center push is
\begin{equation}
 R(X,Y)=(-\lambda^2X,Y+2\lambda X).
 \label{eq:direct-three-vertex-root-map}
\end{equation}
Two consecutive leaf pushes commute, restore equality of the leaves, and
have the combined map
\begin{equation}
 L(X,Y)=(X+2\lambda Y,-\lambda^2Y).
 \label{eq:direct-three-vertex-leaf-map}
\end{equation}

We first justify that the literal coordinate order realizes every such $L$.
Immediately before a leaf pair, suppose the center value is $X<0$, the
common leaf value is $Y>0$, and put $u=Y/(-X)$.  After pushing one leaf, the
untouched leaf remains $Y$, the pushed leaf is $-\lambda^2Y$, and the center
is $X+\lambda Y$.  The pushed leaf has smaller rank than the untouched one,
and
\begin{equation}
 \frac{|X+\lambda Y|}{Y}
 =\left|\lambda-\frac1u\right|.
 \label{eq:direct-three-vertex-between-leaves}
\end{equation}
At the first leaf pair $u=2/\lambda$, so this ratio is $\lambda/2<3/4$.
At the second initial pair, direct multiplication gives
$u=4/(3\lambda)$; at every subsequent pair the crossing bounds below give
$4/3<u<1.362$.  In both cases the ratio is less than $0.267<3/4$.  Since the center and
leaf rank weights are $2/3$ and $1/2$, respectively, the untouched leaf
strictly outranks both other vertices and is the next singleton batch.
Thus \cref{eq:direct-three-vertex-leaf-map} is the exact literal map.

Starting from $(X,Y)=(1,0)$, the first two restored cycles are
$R,L,R,L$.  Direct multiplication gives
\begin{equation}
 (X,Y)=(5\lambda^4,-4\lambda^5),
 \qquad t:=\frac YX=-\frac{4\lambda}{5}.
 \label{eq:direct-three-vertex-entry-state}
\end{equation}
For $h\leq10^{-2}$, $t$ belongs to $I=[-0.82,-0.77]$ and $X>0$.

Consider any restored state with $X>0$ and $t\in I$.  After $k$ consecutive
center pushes,
\begin{equation*}
 X_k=(-\lambda^2)^kX,
 \qquad Y_k=Y+e[1-(-\lambda^2)^k]X.
\end{equation*}
Write $a_k=\lambda^{2k}$.  The live ratio is
\begin{equation}
 \frac{Y_k}{X_k}
 =\begin{cases}
  -e-(t+e)/a_k,&k\text{ odd},\\
  -e+(t+e)/a_k,&k\text{ even}.
 \end{cases}
 \label{eq:direct-three-vertex-root-run-ratio}
\end{equation}
The root is ranked first exactly while the absolute ratio is at most
$c:=4/3$.  Define $A(t):=(t+e)/(c-e)$.  For $\lambda\geq0.98$ and
$t\in I$, direct interval arithmetic gives
\begin{equation}
 0.53<A(t)<0.70.
 \label{eq:direct-three-vertex-A-interval}
\end{equation}
Every even ratio before the crossing has magnitude below $0.83<c$; the odd
ratios increase in magnitude as $a_k$ decreases.  Hence the first leaf wins
after the first odd integer $m$ satisfying $a_m<A(t)$.  The preceding odd
iterate gives
\begin{equation}
 \lambda^4A(t)\leq a_m<A(t),
 \qquad 0.48<a_m<0.70.
 \label{eq:direct-three-vertex-crossing-window}
\end{equation}
At the crossing write $Y_m=u a_mX$ and $X_m=-a_mX$.  Then
\begin{equation}
 \frac43<u\leq e+(c-e)\lambda^{-4}<1.362.
 \label{eq:direct-three-vertex-u-window}
\end{equation}
The two singleton leaf batches therefore execute consecutively.  Their
restored output ratio is
\begin{equation}
 t^+=-\frac{\lambda^2u}{2\lambda u-1}
 \in[-0.800,-0.783]\subset I.
 \label{eq:direct-three-vertex-invariant-interval}
\end{equation}
Thus $I$ is invariant and every subsequent macrocycle is exactly $LR^m$
with odd $m$.  Moreover, $a_m<0.70$ and
$\log(1/\lambda)=2\operatorname{arctanh}(h)\leq2h/(1-h^2)$ imply
\begin{equation}
 m>\frac{\log(10/7)}{2\log(1/\lambda)}
 \geq\frac1{12h}.
 \label{eq:direct-three-vertex-root-run-lower}
\end{equation}

It remains to show that many macrocycles occur above the gate.  At restored
states put $S=X+Y$.  A center push at value $V$ changes $S$ by
$-(1-\lambda)^2V$.  Each individual leaf push changes
$X+(Y_1+Y_2)/2$ by one half of the same expression, so the two-leaf block
has the same identity with common value $Y$.  Put
$\delta=(1-\lambda)^2$.  Over one odd center run,
\begin{equation*}
 \sum_{j=0}^{m-1}X_j=X\frac{1+a_m}{1+\lambda^2}>0,
 \qquad Y_m=u a_mX>0.
\end{equation*}
Consequently one macrocycle satisfies
\begin{equation}
 S^+=S-\delta X\left[
  \frac{1+a_m}{1+\lambda^2}+u a_m
 \right]
 \geq(1-11\delta)S.
 \label{eq:direct-three-vertex-slow-sum}
\end{equation}
Indeed $t\in I$ gives $S=X(1+t)>0$ and $X\leq S/0.18$, while the crossing
bounds control the bracket.  Every partial center run has a positive
geometric partial sum at most $2X/(1+\lambda^2)$, so within a macrocycle
$S_k\geq(1-6\delta)S$.

After \cref{eq:direct-three-vertex-entry-state},
$S_0=\lambda^4(5-4\lambda)\geq0.92$.  Let
$M=\lfloor1/(44\delta)\rfloor$.  Bernoulli's inequality gives, at the first
$M$ restored boundaries,
\begin{equation*}
 S_j\geq(1-11\delta)^jS_0\geq\frac34S_0,
\end{equation*}
and the partial-run bound keeps $S>0.68$ between them.  If all three
coordinates were below the spreading gate, then
$|S|\leq|X|+(|Y_1|+|Y_2|)/2\leq2\beta=0.1$, a contradiction.  Thus all $M$
macrocycles execute and every selected coordinate is live.  Since
$\delta=(2h/(1+h))^2\leq4h^2$ and $h\leq10^{-2}$,
\begin{equation*}
 M\geq\frac1{352h^2}.
\end{equation*}
Each macrocycle has at least $1/(12h)$ center pushes of charge three.
Therefore
\begin{equation*}
 \Work_{\rm spread}
 \geq\frac{3}{12h}\frac1{352h^2}
 =\frac1{1408h^3},
\end{equation*}
which proves \cref{eq:direct-three-vertex-lower-bound}.
\end{proof}

For odd $m$ and $a=\lambda^{2m}$, the exact macro matrix is
\begin{equation}
 K_m:=LR^m
 =\begin{pmatrix}
 -a+\dfrac{4\lambda^2(1+a)}{1+\lambda^2}&2\lambda\\[5pt]
 -\dfrac{2\lambda^3(1+a)}{1+\lambda^2}&-\lambda^2
 \end{pmatrix},
 \label{eq:direct-three-vertex-macro-matrix}
\end{equation}
with
\begin{equation}
 \det K_m=a\lambda^2,
 \qquad
 \det(I-K_m)=
 \frac{(1+a)(1-\lambda^2)^2}{1+\lambda^2}.
 \label{eq:direct-three-vertex-macro-spectrum}
\end{equation}
The greedy crossing keeps $a\in(0.48,0.70)$.  As $\alpha\downarrow0$ the
macro eigenvalues tend to $1$ and $a$: the slow mode contracts only by
$\Theta(\alpha)$ while a macrocycle costs
$\Theta(1/\sqrt\alpha)$ center pushes.  The slow-sum proof remains valid
when the odd run length changes between cycles.

The positive results and this lower bound now give a sharp structural split.
Uniform expansion is accelerated by a final-region gap; the gapless endpoint
path is accelerated by its causal event wave; but a reached reflecting leaf
can combine $\Theta(1/\sqrt\alpha)$ greedy echoes with a
$1-\Theta(\alpha)$ slow macro mode.
